\documentclass[11pt,a4paper]{article}
\usepackage[utf8]{inputenc}
\usepackage[spanish]{babel}
\usepackage[T1]{fontenc}
\usepackage{geometry}
\geometry{margin=2cm}
\usepackage{xcolor}
\usepackage{listings}
\usepackage{tcolorbox}
\tcbuselibrary{most}
\usepackage{booktabs}
\usepackage{array}
\usepackage{hyperref}
\usepackage{fancyhdr}
\usepackage{titlesec}
\usepackage{multicol}

\definecolor{hookpurple}{RGB}{90,50,150}
\definecolor{codebg}{RGB}{245,245,245}
\definecolor{alertred}{RGB}{220,50,50}
\definecolor{tipgreen}{RGB}{39,174,96}
\definecolor{warnambel}{RGB}{230,126,34}
\definecolor{darkgray}{RGB}{50,50,50}

\lstdefinestyle{python}{
  language=Python,
  backgroundcolor=\color{codebg},
  basicstyle=\ttfamily\small,
  keywordstyle=\color{hookpurple}\bfseries,
  stringstyle=\color{teal},
  commentstyle=\color{gray}\itshape,
  numbers=left, numberstyle=\tiny\color{gray},
  frame=single, framerule=0.4pt, rulecolor=\color{hookpurple!40},
  breaklines=true, tabsize=2, showstringspaces=false
}

\lstdefinestyle{shell}{
  language=bash,
  backgroundcolor=\color{darkgray},
  basicstyle=\ttfamily\small\color{white},
  keywordstyle=\color{yellow}\bfseries,
  numbers=none,
  frame=none,
  breaklines=true, tabsize=2
}

% Soporte UTF-8 en listings (español)
\lstset{
  literate=
    {á}{{\'{a}}}1 {é}{{\'{e}}}1 {í}{{\'{i}}}1 {ó}{{\'{o}}}1 {ú}{{\'{u}}}1
    {Á}{{\'{A}}}1 {É}{{\'{E}}}1 {Í}{{\'{I}}}1 {Ó}{{\'{O}}}1 {Ú}{{\'{U}}}1
    {ñ}{{\~{n}}}1 {Ñ}{{\~{N}}}1 {ü}{{\"u}}1
    {¿}{{?`}}1 {¡}{{!`}}1
}

\newtcolorbox{tipbox}{colback=tipgreen!10,colframe=tipgreen,title=\textbf{Tip},fonttitle=\small\bfseries}
\newtcolorbox{alertbox}{colback=alertred!10,colframe=alertred,title=\textbf{Crítico},fonttitle=\small\bfseries}
\newtcolorbox{ejerciciobox}[1]{colback=warnambel!8,colframe=warnambel,title=\textbf{Ejercicio #1},fonttitle=\small\bfseries}
\newtcolorbox{solucionbox}[1]{colback=tipgreen!8,colframe=tipgreen!60,title=\textbf{Solución #1},fonttitle=\small\bfseries}

\pagestyle{fancy}
\fancyhf{}
\lhead{\textcolor{hookpurple}{\textbf{Webhooks e Integraciones}}}
\rhead{\textcolor{gray}{Preparación Técnica — Toku}}
\cfoot{\thepage}

\titleformat{\section}{\large\bfseries\color{hookpurple}}{}{0em}{}[\titlerule]
\titleformat{\subsection}{\normalsize\bfseries\color{darkgray}}{}{0em}{}

\begin{document}

\begin{titlepage}
  \centering
  \vspace*{2cm}
  {\Huge\bfseries\color{hookpurple} Webhooks e Integraciones\par}
  \vspace{0.5cm}
  {\large\color{gray} Diseño robusto de sistemas de notificación y APIs\par}
  \vspace{0.3cm}
  \textcolor{hookpurple!40}{\rule{8cm}{2pt}}
  \vspace{1.5cm}

  \begin{tcolorbox}[colback=hookpurple!5,colframe=hookpurple,width=11cm]
    \centering
    \textbf{Foco:} Lo que el rol de Integrations Engineer exige\\[4pt]
    HMAC, idempotencia, retries, async processing, OAuth 2.0
  \end{tcolorbox}

  \vspace{1.5cm}
  \begin{tabular}{ll}
    \textcolor{hookpurple}{\textbullet} & Diseño de webhooks robustos\\
    \textcolor{hookpurple}{\textbullet} & Verificación HMAC-SHA256\\
    \textcolor{hookpurple}{\textbullet} & Procesamiento asíncrono con colas\\
    \textcolor{hookpurple}{\textbullet} & OAuth 2.0 y autenticación de APIs\\
    \textcolor{hookpurple}{\textbullet} & Manejo de reintentos y backoff\\
    \textcolor{hookpurple}{\textbullet} & Integración con ERPs y sistemas externos\\
  \end{tabular}
  \vfill
  {\small\color{gray} Marzo 2026 — Estudio Personal\par}
\end{titlepage}

\tableofcontents
\newpage

% ─── SECCIÓN 1: WEBHOOKS ─────────────────────────────────────────────────────
\section{Webhooks: Diseño Robusto}

\begin{tcolorbox}[colback=hookpurple!5,colframe=hookpurple]
Un webhook es un \textbf{HTTP POST que un sistema envía a tu servidor} cuando algo ocurre.
Es la diferencia entre \textit{polling} (preguntar "¿pasó algo?") y \textit{push} (te aviso cuando pasa).
Toku usa webhooks para notificarte de pagos, errores, nuevos payment methods, etc.
\end{tcolorbox}

\subsection{Las 5 Reglas de Oro de los Webhooks}

\begin{enumerate}
  \item \textbf{Responde 200 inmediatamente} — no proceses en el handler, solo acepta y encola
  \item \textbf{Verifica la firma siempre} — cualquier request puede llegar a tu endpoint
  \item \textbf{Sé idempotente} — el mismo evento puede llegar más de una vez
  \item \textbf{Loggea todo} — en pagos, el debugging depende de los logs
  \item \textbf{Usa HTTPS} — nunca HTTP para webhooks de pagos
\end{enumerate}

\subsection{Implementación completa con FastAPI}

\begin{lstlisting}[style=python]
# webhooks/router.py
from fastapi import APIRouter, Request, HTTPException, BackgroundTasks
from fastapi.responses import JSONResponse
import hmac
import hashlib
import logging
from datetime import datetime

from database import db
from tasks import procesar_evento_pago

router = APIRouter(prefix="/webhooks", tags=["webhooks"])
logger = logging.getLogger(__name__)

WEBHOOK_SECRET = "tu_secreto_toku"

def verificar_firma(body: bytes, signature: str) -> bool:
    """HMAC-SHA256: verificar que el request viene de Toku"""
    expected = hmac.new(
        WEBHOOK_SECRET.encode("utf-8"),
        body,
        hashlib.sha256
    ).hexdigest()
    # compare_digest evita timing attacks
    return hmac.compare_digest(expected, signature)

@router.post("/toku")
async def recibir_webhook(
    request: Request,
    background_tasks: BackgroundTasks
):
    # 1. Obtener body raw (importante: ANTES de parsear JSON)
    body = await request.body()
    signature = request.headers.get("X-Toku-Signature", "")

    # 2. Verificar firma
    if not verificar_firma(body, signature):
        logger.warning("Firma invalida recibida en webhook")
        raise HTTPException(status_code=401, detail="Firma invalida")

    # 3. Parsear evento
    evento = await request.json()
    event_id = evento.get("id")
    event_type = evento.get("type")

    logger.info(f"Webhook recibido: {event_type} | {event_id}")

    # 4. Idempotencia: registrar el event_id
    insertado = await db.execute(
        """INSERT INTO webhook_events (event_id, event_type, payload)
           VALUES ($1, $2, $3)
           ON CONFLICT (event_id) DO NOTHING""",
        event_id, event_type, evento
    )

    if insertado == "INSERT 0 0":
        # Duplicado: ya fue procesado
        logger.info(f"Evento duplicado ignorado: {event_id}")
        return JSONResponse({"status": "duplicate"})

    # 5. Encolar procesamiento asíncrono (background task)
    background_tasks.add_task(procesar_evento_pago, evento)

    # 6. Responder 200 inmediatamente (no esperar procesamiento)
    return JSONResponse({"status": "received"})
\end{lstlisting}

\begin{lstlisting}[style=python]
# tasks/payment_tasks.py
from database import db
import logging

logger = logging.getLogger(__name__)

async def procesar_evento_pago(evento: dict):
    """Procesamiento real del evento (fuera del request HTTP)"""
    event_type = evento["type"]
    data = evento["data"]

    try:
        if event_type == "invoice.paid":
            await manejar_pago_exitoso(data)
        elif event_type == "invoice.failed":
            await manejar_pago_fallido(data)
        elif event_type == "payment_method.created":
            await manejar_nuevo_metodo_pago(data)
        elif event_type == "transaction.completed":
            await manejar_transaccion_completada(data)

        # Marcar como procesado
        await db.execute(
            "UPDATE webhook_events SET estado='completado', procesado_at=NOW() WHERE event_id=$1",
            evento["id"]
        )
    except Exception as e:
        logger.error(f"Error procesando evento {evento['id']}: {e}")
        await db.execute(
            "UPDATE webhook_events SET estado='error', error_detalle=$1 WHERE event_id=$2",
            str(e), evento["id"]
        )
        raise  # Para que el retry system lo capture

async def manejar_pago_exitoso(data: dict):
    invoice_id = data["invoice_id"]
    transaction_id = data["transaction_id"]
    monto = data["amount"]

    async with db.transaction():
        # Actualizar factura
        rows = await db.execute(
            """UPDATE invoices
               SET estado='paid', pagada_at=NOW(), updated_at=NOW()
               WHERE id=$1 AND estado='pending'""",
            invoice_id
        )
        if rows == "UPDATE 0":
            logger.warning(f"Invoice {invoice_id} no estaba pending")
            return  # Idempotente: si ya fue pagada, OK

        # Registrar el pago en tu sistema
        await db.execute(
            "INSERT INTO pagos_confirmados (invoice_id, transaction_id, monto) VALUES ($1,$2,$3)",
            invoice_id, transaction_id, monto
        )

    logger.info(f"Pago exitoso procesado: {invoice_id}")
\end{lstlisting}

% ─── SECCIÓN 2: AUTENTICACIÓN DE APIs ────────────────────────────────────────
\section{Autenticación de APIs Externas}

\subsection{API Keys (el más simple)}

\begin{lstlisting}[style=python]
import httpx
import os

# NUNCA hardcodear
TOKU_API_KEY = os.environ["TOKU_API_KEY"]

async def llamar_toku_api(endpoint: str, payload: dict) -> dict:
    async with httpx.AsyncClient() as client:
        response = await client.post(
            f"https://api.trytoku.com{endpoint}",
            headers={
                "Authorization": f"Bearer {TOKU_API_KEY}",
                "Content-Type": "application/json",
            },
            json=payload,
            timeout=30.0
        )
        response.raise_for_status()
        return response.json()
\end{lstlisting}

\subsection{OAuth 2.0 — El Estándar de Integraciones}

\begin{tcolorbox}[colback=hookpurple!5,colframe=hookpurple,title=\textbf{¿Cuándo aparece OAuth 2.0 en Toku?}]
Cuando Toku se integra con el sistema del cliente (ERP, CRM),
a veces ese sistema externo requiere OAuth 2.0 para autorizar el acceso.
El rol de Integrations Engineer tiene que manejar esto.
\end{tcolorbox}

\begin{lstlisting}[style=python]
# Flujo Client Credentials (machine-to-machine, el mas comun en integraciones)
import httpx
import time

class OAuthClient:
    def __init__(self, client_id: str, client_secret: str, token_url: str):
        self.client_id = client_id
        self.client_secret = client_secret
        self.token_url = token_url
        self._token = None
        self._token_expiry = 0

    async def get_token(self) -> str:
        """Obtener token, renovando si expiró"""
        now = time.time()
        if self._token and now < self._token_expiry - 60:  # 60s buffer
            return self._token

        async with httpx.AsyncClient() as client:
            response = await client.post(
                self.token_url,
                data={
                    "grant_type": "client_credentials",
                    "client_id": self.client_id,
                    "client_secret": self.client_secret,
                    "scope": "read:pagos write:pagos"
                }
            )
            response.raise_for_status()
            data = response.json()

        self._token = data["access_token"]
        self._token_expiry = now + data["expires_in"]
        return self._token

    async def get(self, url: str) -> dict:
        token = await self.get_token()
        async with httpx.AsyncClient() as client:
            response = await client.get(
                url,
                headers={"Authorization": f"Bearer {token}"}
            )
            response.raise_for_status()
            return response.json()
\end{lstlisting}

\subsection{JWT — Verificación de Tokens}

\begin{lstlisting}[style=python]
import jwt
from datetime import datetime, timedelta

SECRET_KEY = "tu_secreto_jwt"
ALGORITHM = "HS256"

def crear_token(user_id: str, expiry_hours: int = 24) -> str:
    payload = {
        "sub": user_id,
        "iat": datetime.utcnow(),
        "exp": datetime.utcnow() + timedelta(hours=expiry_hours)
    }
    return jwt.encode(payload, SECRET_KEY, algorithm=ALGORITHM)

def verificar_token(token: str) -> dict:
    try:
        payload = jwt.decode(token, SECRET_KEY, algorithms=[ALGORITHM])
        return payload
    except jwt.ExpiredSignatureError:
        raise ValueError("Token expirado")
    except jwt.InvalidTokenError:
        raise ValueError("Token invalido")
\end{lstlisting}

% ─── SECCIÓN 3: REINTENTOS Y RESILIENCIA ─────────────────────────────────────
\section{Reintentos, Backoff y Resiliencia}

\subsection{Backoff Exponencial con Jitter}

\begin{lstlisting}[style=python]
import asyncio
import random
from typing import Callable, TypeVar, Any

T = TypeVar("T")

async def con_reintentos(
    fn: Callable[..., Any],
    max_intentos: int = 5,
    base_delay: float = 1.0,
    max_delay: float = 60.0,
    *args,
    **kwargs
):
    """
    Reintenta una funcion con backoff exponencial + jitter.
    Jitter = aleatoriedad para evitar thundering herd
    (todos los clientes reintentando al mismo tiempo)
    """
    for intento in range(max_intentos):
        try:
            return await fn(*args, **kwargs)
        except Exception as e:
            if intento == max_intentos - 1:
                raise  # Ultimo intento: propagar el error

            # Backoff: 1s, 2s, 4s, 8s, 16s...
            delay = min(base_delay * (2 ** intento), max_delay)
            # Jitter: +-20% del delay para desincronizar clientes
            jitter = delay * random.uniform(-0.2, 0.2)
            espera = delay + jitter

            print(f"Intento {intento+1} fallido. Reintentando en {espera:.1f}s")
            await asyncio.sleep(espera)

# Uso:
resultado = await con_reintentos(
    llamar_toku_api,
    max_intentos=3,
    args=("/transactions",),
    kwargs={"payload": {...}}
)
\end{lstlisting}

\subsection{Circuit Breaker — No machucar un servicio caído}

\begin{lstlisting}[style=python]
from enum import Enum
import time

class Estado(Enum):
    CERRADO = "closed"    # Normal: deja pasar requests
    ABIERTO = "open"      # Falla: bloquea todos los requests
    SEMI_ABIERTO = "half_open"  # Probando si se recupero

class CircuitBreaker:
    def __init__(self, umbral_fallas: int = 5, timeout: float = 60.0):
        self.umbral = umbral_fallas
        self.timeout = timeout
        self.fallas = 0
        self.ultimo_fallo = None
        self.estado = Estado.CERRADO

    async def llamar(self, fn, *args, **kwargs):
        if self.estado == Estado.ABIERTO:
            if time.time() - self.ultimo_fallo > self.timeout:
                self.estado = Estado.SEMI_ABIERTO
            else:
                raise Exception("Circuit breaker abierto: servicio no disponible")

        try:
            resultado = await fn(*args, **kwargs)
            self._exito()
            return resultado
        except Exception as e:
            self._fallo()
            raise

    def _exito(self):
        self.fallas = 0
        self.estado = Estado.CERRADO

    def _fallo(self):
        self.fallas += 1
        self.ultimo_fallo = time.time()
        if self.fallas >= self.umbral:
            self.estado = Estado.ABIERTO
\end{lstlisting}

% ─── SECCIÓN 4: INTEGRACIÓN CON ERP/SISTEMAS EXTERNOS ────────────────────────
\section{Integración con ERP y Sistemas Externos}

\begin{tcolorbox}[colback=warnambel!5,colframe=warnambel,title=\textbf{El desafío del rol de Integrations Engineer}]
Los clientes de Toku tienen sus propios sistemas (SAP, Salesforce, sistemas legacy en COBOL,
bases de datos Oracle...). El trabajo es conectar ese sistema con la API de Toku de forma
que el flujo de datos sea confiable, sin duplicados y fácil de operar.
\end{tcolorbox}

\subsection{Patron: Sync bidireccional confiable}

\begin{lstlisting}[style=python]
# Patron: cursor de sincronizacion
# En lugar de sincronizar TODO siempre, solo sincronizar LO NUEVO

async def sync_clientes_desde_erp(erp_client, toku_client, db):
    """Sincronizar clientes nuevos/actualizados desde ERP a Toku"""

    # Obtener ultimo cursor (timestamp de la ultima sync exitosa)
    cursor = await db.fetchval(
        "SELECT MAX(ultima_sync) FROM sync_state WHERE entidad='customers'"
    ) or "2000-01-01T00:00:00Z"

    # Pedir al ERP solo lo que cambio desde el cursor
    clientes_nuevos = await erp_client.get_customers_since(cursor)

    exitosos = 0
    fallidos = 0

    for cliente in clientes_nuevos:
        try:
            # Verificar si ya existe en Toku (por external_id)
            existente = await toku_client.get_customer_by_external_id(
                cliente["id_erp"]
            )

            if existente:
                # Actualizar
                await toku_client.update_customer(existente["id"], {
                    "nombre": cliente["nombre"],
                    "email": cliente["email"]
                })
            else:
                # Crear nuevo
                await toku_client.create_customer({
                    "external_id": cliente["id_erp"],
                    "nombre": cliente["nombre"],
                    "email": cliente["email"],
                    "telefono": cliente.get("telefono")
                })

            exitosos += 1

        except Exception as e:
            # Loggear y continuar (no abortar toda la sync por un error)
            logger.error(f"Error sincronizando cliente {cliente['id_erp']}: {e}")
            fallidos += 1

    # Actualizar cursor solo si terminamos bien
    await db.execute(
        """INSERT INTO sync_state (entidad, ultima_sync)
           VALUES ('customers', NOW())
           ON CONFLICT (entidad) DO UPDATE SET ultima_sync = NOW()"""
    )

    logger.info(f"Sync completada: {exitosos} ok, {fallidos} fallidos")
\end{lstlisting}

\subsection{Trade-off: API REST vs SFTP vs Low-Code}

\begin{center}
\begin{tabular}{p{3cm}p{4cm}p{4cm}p{3cm}}
\toprule
\textbf{Método} & \textbf{Cuándo usar} & \textbf{Ventajas} & \textbf{Desventajas} \\
\midrule
API REST custom & Tiempo real, sistema moderno & Control total, inmediato & Más desarrollo \\
SFTP + CSV & Batch diario, sistema legacy & Simple, sin cambios en ERP & Latencia de 24h \\
Workato/Tray & Integraciones estándar (Salesforce, SAP) & Rápido de implementar & Costo, límite de complejidad \\
Webhook listener & Reaccionar a eventos de Toku & Eficiente, push & Tu servidor debe estar disponible \\
\bottomrule
\end{tabular}
\end{center}

\section{Cheatsheet: Lo Esencial en una Página}

\begin{multicols}{2}
\textbf{\textcolor{hookpurple}{HMAC-SHA256}}
\begin{lstlisting}[style=python,numbers=none]
import hmac, hashlib
expected = hmac.new(
    secret.encode(),
    raw_body,
    hashlib.sha256
).hexdigest()
hmac.compare_digest(expected, sig)
\end{lstlisting}

\textbf{\textcolor{hookpurple}{Idempotencia en webhook}}
\begin{lstlisting}[style=python,numbers=none]
INSERT INTO events (event_id, ...)
ON CONFLICT (event_id) DO NOTHING;
-- rowcount == 0 -> duplicado
\end{lstlisting}

\textbf{\textcolor{hookpurple}{Backoff exponencial}}
\begin{lstlisting}[style=python,numbers=none]
delay = min(base * 2**intento, max_d)
jitter = delay * random.uniform(-.2,.2)
await asyncio.sleep(delay + jitter)
\end{lstlisting}

\columnbreak

\textbf{\textcolor{hookpurple}{OAuth2 Client Credentials}}
\begin{lstlisting}[style=python,numbers=none]
POST /token
  grant_type=client_credentials
  client_id=...
  client_secret=...
# -> access_token, expires_in
\end{lstlisting}

\textbf{\textcolor{hookpurple}{Eventos Toku clave}}
\begin{lstlisting}[style=python,numbers=none]
invoice.paid       -> marcar pagado
invoice.failed     -> notificar, retry
payment_method.created -> cobro futuro
transaction.completed  -> confirmacion
settlement.created     -> fondos listos
\end{lstlisting}
\end{multicols}

\end{document}
